% Gunakan setelah ujian dan revisi benar-benar disahkan.
% Lengkapi nama penguji, identitas pejabat, tanggal, serta tanda tangan aktual.
\skripsifrontchapter{PENGESAHAN}
\begin{center}
\textbf{\judulskripsi}\par
\namamahasiswa\par NIM: \nimmahasiswa
\end{center}
Naskah ini dipertahankan di hadapan penguji pada tanggal \rule{3cm}{0.4pt}. Hasil ujian dan penyelesaian revisi diisi sesuai berita acara yang telah disahkan.

\begin{center}
\begin{tabular}{p{5.8cm}p{5.8cm}}
Pembimbing 1 & Pembimbing 2 \\
\rule{0pt}{2cm} & \\
\pembimbingsatu & \pembimbingdua \\
\end{tabular}
\end{center}
\vspace{0.8cm}
\noindent Penguji: \rule{7cm}{0.4pt}\par
\noindent Ketua Departemen: \rule{7cm}{0.4pt}\par
\noindent Tanggal pengesahan: \rule{5cm}{0.4pt}
